\documentclass[12pt,a4paper]{article}

%% PACHETE MATEMATICE
\usepackage{amsmath,amssymb,amsfonts}
\usepackage{amsthm}
\usepackage{mathpazo}
\usepackage{eulervm}
\usepackage{enumerate}
\usepackage{color}
\usepackage{graphicx}
\usepackage[all]{xy}
\usepackage{xcolor}
\usepackage{framed}
\usepackage[unicode]{hyperref}
\setcounter{MaxMatrixCols}{20}
\usepackage{multicol}
% \usepackage[romanian]{babel}


%% ASPECT
\linespread{1.05}
\usepackage[T1]{fontenc}
\usepackage[utf8x]{inputenc}
\usepackage[margin=2cm]{geometry}
% \usepackage[color]{showkeys}
\usepackage{anyfontsize}
\usepackage{titlesec}
\titleformat*{\section}{\large\bfseries}

\usepackage{fancyhdr}
\pagestyle{fancy}
\rhead{\small \color{gray} Notițe de Adrian Manea}
\lhead{\small \color{gray} ETTI-AM2, 2017---2018, M. Joița \& A. Niță}


\usepackage[autostyle = false, style = german]{csquotes}
\usepackage[romanian]{babel}
\newcommand\qq{\enquote}

%% COMENZILE MELE
\renewcommand*{\proofname}{Dem.:}
\newcommand\bb{\mathbb}
\newcommand\dr{\mathrm}
\newcommand\kal{\mathcal}
\newcommand\fk{\mathfrak}
\newcommand\op{\oplus}
\newcommand\ot{\otimes}
\newcommand\ds{\displaystyle}
\newcommand\id{\indent}
\newcommand\nid{\noindent}
\newcommand\seq{\subseteq}
\newcommand\sneq{\subsetneq}
\newcommand\speq{\supseteq}
\newcommand\spneq{\supsetneq}
\newcommand\rar{\rightarrow}
\newcommand\Rar{\Rightarrow}
\newcommand\xrar{\xrightarrow}
\newcommand\lar{\leftarrow}
\newcommand\Lar{\Leftarrow}
\newcommand\xlar{\xleftarrow}
\newcommand\lrar{\leftrightarrow}
\newcommand\Lrar{\Leftrightarrow}
\newcommand\ideal{\trianglelefteq}
\newcommand\ai{\text{ a.\^{i}. }}
\newcommand\sk{(\textit{Schi\c{t}\u{a}}) }
\newcommand\rhup{\rightharpoonup}
\newcommand\rhdn{\rightharpoondown}
\newcommand\lhup{\leftharpoonup}
\newcommand\lhdn{\leftharpoondown}
\newcommand\vid{\emptyset}
\newcommand\wh{\widehat}
\newcommand\ol{\overline}
\newcommand\NN{\mathbb{N}}
\newcommand\RR{\mathbb{R}}
\newcommand\ZZ{\mathbb{Z}}
\newcommand\QQ{\mathbb{Q}}
\newcommand\CC{\mathbb{C}}

%% STILURILE MELE
\newtheoremstyle{dotless}{}{}{\itshape}{}{\bfseries}{:}{ }{}

\theoremstyle{dotless}
\newtheorem{theorem}{Teorem\u{a}}
\newtheorem{proposition}{Propozi\c{t}ie}
\newtheorem{lemma}{Lem\u{a}}
\newtheorem{corollary}{Corolar}
\newtheorem{exercise}{Exerci\c{t}iu}

\newtheoremstyle{dotlessdr}{}{}{}{}{\bfseries}{:}{ }{}
\theoremstyle{dotlessdr}
\newtheorem{example}{Exemplu}
\newtheorem{definition}{Defini\c{t}ie}
\newtheorem{remark}{Observa\c{t}ie}

\newenvironment{solution}
{\begin{proof}[Soluție:\nopunct]}
  {\end{proof}}


\begin{document}
% \definecolor{labelkey}{RGB}{222,0,0}

\begin{center}
  {\large\textbf{Seminar 10 \\
      Transformarea Fourier}}
\end{center}

\vspace{2cm}

\section*{Definiții}

\indent\indent Începem studiul \emph{transformărilor integrale}, folosind analiza
complexă, cu transformările Fourier, utile, ca și seriile Fourier, pentru studiul
funcțiilor periodice, de tipul semnalelor electrice.

Facem următoarea notație a unei mulțimi de funcții pe care o vom folosi în continuare:
\[
  L^1( \RR ) = \Big\{ f : \RR \to \CC \Big\mid \int_{-\infty}^\infty |f(t)| dt < \infty \Big\}.
\]

Cu aceasta, avem:

\begin{definition}\label{def:tr-four}
  Se numește \emph{transformarea Fourier} a funcției $ f \in L^1(\RR) $ o funcție
  complexă, notată cu $ \kal{F}[f] : \RR \to \CC $, definită prin:
  \begin{equation}
    \label{eq:transf-f}
    \kal{F}[f](\omega) = \int_{-\infty}^\infty f(t) e^{i\omega t} dt.
  \end{equation}

  Funcția $ \kal{F}[f](\omega) $ se mai numește \emph{funcția spectrală} sau
  \emph{spectrul (în frecvență)} asociat(ă) funcției $ f $. Dacă privim $ f $ ca
  un semnal, dependent de timp, transformarea Fourier îi asociază spectrul acestuia.

  O notație alternativă pentru $ \kal{F}[f](\omega) $ este $ \widehat{f}(\omega) $.
\end{definition}

Ideea de bază a unei transformări Fourier este următoarea. Dacă se dă o funcție periodică
(sau convertită la una periodică printr-un artificiu de repetiție, practic), ei i se
asociază \emph{seria Fourier}. Aceasta o aproximează cu o serie de sinusuri și cosinusuri,
funcțiile periodice cel mai des întîlnite. Practic, are loc o superpoziție de termeni cu
sinusuri și cosinusuri, un fel de interferență a undelor electromagnetice. La pasul
următor, \emph{transformata Fourier} preia minimele și maximele acestor \qq{interferențe},
iar rezultatul este un semnal (aproape) discret (\qq{digitalizat}), care reprezintă
valorile cele mai importante din semnal.\footnote{O explicație animată este dată \href{https://en.wikipedia.org/wiki/Fourier_transform\#/media/File:Fourier_transform_time_and_frequency_domains_(small).gif}{aici}.}

Conform teoriei seriilor Fourier, dacă funcția dată este pară, definiția se reduce
la cazul mai simplu al cosinusurilor:
\begin{equation}
  \label{eq:transf-f-cos}
  \widehat{f}(\omega) = 2 \int_0^\infty f(t) \cos \omega t dt, \quad \omega \in \RR.
\end{equation}

Similar, dacă funcția este impară, dezvoltarea conține doar funcții sinus:
\begin{equation}
  \label{eq:transf-f-sin}
  \widehat{f}(\omega) = -2i \int_0^\infty f(t) \sin \omega t dt, \quad \omega \in \RR.
\end{equation}

Putem recupera funcția dată din transformata Fourier printr-o \emph{formulă de inversare}:

\begin{theorem}[Formula Fourier de inversare]\label{thm:fourier-inv}
  Fie $ f : \RR \to \RR $ o funcție din $ L^1(\RR) $. Dacă $ \widehat{f}(\omega) $ este
  transformata ei Fourier și presupunînd că $ \widehat{f}(\omega) \in L^1(\CC) $,
  atunci funcția inițială se obține din:
  \begin{equation}
    \label{eq:transf-f-inv}
    f(t) = \frac{1}{2\pi} \int_{-\infty}^\infty \widehat{f}(\omega) e^{it\omega} d\omega, %
    \quad \forall t \in \RR
  \end{equation}
\end{theorem}

În particular, rezultă că putem calcula și transformata Fourier prin inversarea
inversei:
\begin{equation}
  \label{eq:transf-f-iinv}
  \widehat{f}(\omega) = \frac{1}{\sqrt{2\pi}}
\end{equation}

%%%%%%%%%%%%%%%%%%%%%%%%%%%%%%%%%%%%%%%%%%%%%%%%%%%%%%%%%%%%%%%%%%%%%%

\section*{Proprietăți}

\indent\indent Următoarele proprietăți sînt esențiale pentru transformarea Fourier.
Dacă nu se precizează altfel, vom nota prin $ \widehat{f}(\omega) $ transformata
Fourier a funcției $ f $ (și similar pentru orice altă funcție), iar relația dintre
cele două se va nota cu $ f(t) \leftrightarrow \widehat{f}(\omega) $ (și similar
pentru orice altă funcție):

\begin{itemize}
\item \textbf{Liniaritatea}: Dacă $ \alpha, \beta \in \CC $, iar $ f, g \in L^1(\RR) $,
  atunci:
  \[
    \alpha f(t) + \beta g(t) \leftrightarrow \alpha \widehat{f}(\omega) + %
    \beta \widehat{g}(\omega);
  \]
\item \textbf{Simetria:} $ \widehat{f}(t) \leftrightarrow 2 \pi f(-\omega) $;
\item \textbf{Schimbarea de scală}: Dacă $ \alpha \in \CC $, atunci
  \[
    f(\alpha t) \leftrightarrow \frac{1}{|\alpha|} \widehat{f} \Big(\frac{\omega}{\alpha}\Big);
  \]
\item \textbf{Translația timpului}: Dacă $ t_0 \in \RR $, atunci avem:
  \[
    f(t - t_0) \leftrightarrow \widehat{f}(\omega) e^{-it_0\omega};
  \]
\item \textbf{Translația frecvenței}: Dacă $ \omega_0 \in \RR $, atunci:
  \[
    e^{i\omega_0 t} f(t) \leftrightarrow \widehat{f}(\omega - \omega_0);
  \]
\item \textbf{Derivarea în raport cu timpul}: Dacă $ f $ este de $ n $ ori derivabilă,
  atunci:
  \[
    \frac{d^n f}{dt^n} \leftrightarrow (i\omega)^m \widehat{f}(\omega);
  \]
\item \textbf{Derivarea în raport cu frecvența}: Dacă funcțiile $ f(t), tf(t), \dots, %
  t^n f(t) $ sînt integrabile pe $ \RR $, atunci:
  \[
    (-it)^n f(t) \leftrightarrow \frac{d^n \widehat{f}(\omega)}{d\omega^n};
  \]
\item \textbf{Transformarea conjugatei complexe}: Fie $ f^\ast(t) $ conjugata
  complexă a funcției $ f $. Atunci are loc:
  \[
    f^\ast(t) \leftrightarrow \widehat{f}^\ast(-\omega);
  \]
\item \textbf{Teorema de convoluție în timp}: Definim \emph{produsul de convoluție}
  al funcțiilor $ f $ și $ g $ prin:
  \[
    h(t) = (f \ast g)(t) = \int_{-\infty}^\infty f(\tau)g(t - \tau) d\tau.
  \]
  Atunci are loc:
  \[
    (f \ast g)(t) \leftrightarrow \widehat{f}(\omega) \cdot \widehat{g}(\omega).
  \]
\item \textbf{Teorema de convoluție în frecvență}: În condițiile proprietății precedente,
  are loc:
  \[
    f(t) g(t) \leftrightarrow %
    \frac{1}{2\pi} \int_{-\infty}^\infty \widehat{f}(y) \widehat{g}(\omega - y) dy;
  \]
\item $ \widehat{f} $ este uniform continuă pe $ \RR $ și în plus, are loc:
  \[
    \lim_{|\omega| \to \infty} |\widehat{f}(\omega)| = 0;
  \]
\item Dacă $ f_n : \RR \to \RR, n \in \NN^\ast $ este un șir de funcții convergent
  către funcția $ f : \RR \to \RR $ în spațiul $ L^1(\RR) $, adică:
  \[
    \lim_{n \to \infty} \int_{-\infty}^\infty |f_n(t) - f(t)| dt = 0,
  \]
  atunci șirul $ \widehat{f}_n $ converge uniform către $ \widehat{f} $, adică:
  \[
    \lim_{n \to \infty} \sup_{\omega \in \RR} |\widehat{f}_n(\omega) - \widehat{f}(\omega)| = 0.
  \]
\end{itemize}

%%%%%%%%%%%%%%%%%%%%%%%%%%%%%%%%%%%%%%%%%%%%%%%%%%%%%%%%%%%%%%%%%%%%%%

\section*{Exemple rezolvate}

1. Să se calculeze transformata Fourier a funcției $ f(t) = e^{-a|t|}, a > 0 $.

\begin{solution} Conform definiției, avem succesiv:
  \begin{align*}
    \widehat{f}(\omega) &= \int_{-\infty}^\infty e^{-|t|} e^{-i\omega t} dt \\
                        &= \int_{-\infty}^\infty e^{-|t|} (\cos \omega t - i \sin \omega t) dt \\
                        &= 2 \int_0^{\infty} e^{-t} \cos \omega t dt \\
                        &= \int_0^\infty e^{-t} (e^{i\omega t} + e^{-i\omega t} dt \\
                        &= -\frac{1}{i\omega - 1} + \frac{1}{i \omega + 1} \\
                        &= \frac{2}{\omega^2 + 1}.
  \end{align*}

  Folosind formula de schimbare de scală, obținem:
  \[
    \widehat{f}(\omega) = \frac{1}{a} \widehat{f} \Big( \frac{\omega}{a} \Big) = %
    \frac{2}{a \cdot \big( \frac{\omega^2}{a^2} + 1 \big)} = %
    \frac{2a}{\omega^2 + a^2}.
  \]
\end{solution}

\vspace{1cm}

2. Să se calculeze transformata Fourier a funcției $ f(t) = e^{-7|t + 4|} $.

\begin{solution}
  Observăm că putem folosi exemplul anterior, cu $ a = 7 $. De asemenea, conform teoremei
  de translație a timpului, cu $ t_0 = 4 $ și teoremei de schimbare de scală obținem în final:
  \[
    \widehat{f} = \frac{14 e^{4i\omega}}{\omega^2 + 49}.
  \]
\end{solution}

\vspace{1cm}

3. Să se calculeze transformata Fourier a funcției $ f(t) = te^{-\alpha t^2}, \alpha > 0 $.

\begin{solution}
  Pornim cu funcția mai simplă $ f_1(t) = e^{-t^2} $ și calculăm:
  \begin{align*}
    \widehat{f}_1(\omega) &= \int_{-\infty}^\infty e^{-t^2} e^{-i\omega t} dt \\
                          &= \int_{-\infty}^\infty e^{-t^2} \cos \omega t dt -
                            i \int_{-\infty}^\infty e^{-t^2} \sin \omega t dt \\
                          &= 2 \int_0^\infty e^{-t^2} \cos \omega t dt.
  \end{align*}
  Atunci, prin derivare în interiorul integralei (după $ \omega $), găsim:
  \[
    \widehat{f}'_1(\omega) = -2\int_0^\infty e^{-t^2} t \sin \omega t dt = %
    -\frac{1}{2} \omega \widehat{f}_1(\omega),
  \]
  ultima egalitate rezultînd după integrare prin părți.

  Așadar, am ajuns la o ecuație diferențială cu variabile separabile:
  \[
    \frac{d\widehat{f}_1}{d \omega} = -\frac{1}{2} \omega \widehat{f}_1,
  \]
  pe care o rezolvăm și obținem $ \widehat{f}_1(\omega) = ce^{\frac{-\omega^2}{4}} $.

  Constanta se poate calcula ținînd cont de egalitatea:
  \[
    c = \widehat{f}_1(0) = \int_{-\infty}^\infty e^{-t^2} dt = \sqrt{\pi} \Rightarrow %
    \widehat{f}_1(\omega) = \sqrt{\pi} e^{-\omega^2}{4}.
  \]

  Deoarece $ e^{-\alpha t^2} = e^{-( t \sqrt{\alpha})^2} $, putem folosi formula de
  schimbare de scală și obținem:
  \[
    \kal{F}[e^{-\alpha t^2}](\omega) = \frac{1}{\sqrt{\alpha}} %
    \widehat{f}_1\Big( \frac{\omega}{\sqrt{\alpha}} \Big) = %
    \frac{1}{\sqrt{\alpha}} \sqrt{\pi} e^{-\frac{\omega^2}{4 \alpha}}.
  \]

  Acum, cu teorema de derivare în raport cu frecvența, avem:
  \[
    \widehat{f}(\omega) = %
    i \sqrt{\frac{\pi}{\alpha}} \Big(-\frac{\omega}{2 \alpha}\Big) e^{-\omega^2}{4 \alpha} = %
    -\frac{i}{2} \sqrt{\pi}{\alpha^3} \omega e^{-\omega^2}{4 \alpha}.
  \]
\end{solution}

\vspace{1cm}

4. Să se calculeze transformata Fourier a \emph{impulsului triunghiular de lungime $ 2T $}:
\[
  q_T(t) =
  \begin{cases}
    1 - \frac{|t|}{T}, & |t| \leq T \\
    0, & |t| > T
  \end{cases}.
\]

\begin{solution}
  Deoarece $ q_T $ este funcție pară, putem calcula direct folosind formula cu cosinusuri
  (\eqref{eq:transf-f-cos}) și obținem:
  \begin{align*}
    \widehat{f}(\omega) &= 2 \int_0^T \Big( 1 - \frac{t}{T} \Big) \cos \omega t dt \\
                        &= 2 \int_0^T \frac{\sin \omega t}{T \omega} dt \\
                        &= 2 \frac{1}{T \omega} \int_0^T \sin \omega t dt \\
                        &= \frac{2(1 - \cos \omega T)}{T \omega^2} \\
                        &= \frac{4 \sin^2 \frac{T\omega}{2}}{T \omega^2}.
  \end{align*}
\end{solution}

\vspace{1cm}

5. Să se calculeze transformata Fourier a funcției \emph{sinus cardinal}:
\[
  \sin c(t) = \frac{\sin t}{t}.
\]

\begin{solution}
  Funcția este pară, deci putem folosi din nou formula cu seria de cosinusuri
  (\eqref{eq:transf-f-cos}):
  \begin{align*}
    \widehat{f}(\omega) &= 2 \int_0^\infty \frac{\sin t}{t} \cos \omega t dt \\
                        &= \int_0^\infty \frac{\sin t(1 + \omega) + \sin t(1 - \omega)}{t} dt \\
                        &= \frac{\pi}{2} \big[ \mathrm{sgn}(1 + \omega) + \mathrm{sgn}(1 - \omega) \big],
  \end{align*}
  folosind relația:
  \[
    \int_0^\infty \frac{\sin at}{t} dt = \frac{\pi}{2} \mathrm{sgn} a.
  \]
\end{solution}

%%%%%%%%%%%%%%%%%%%%%%%%%%%%%%%%%%%%%%%%%%%%%%%%%%%%%%%%%%%%%%%%%%%%%%

\section*{Tabel de transformate Fourier}

\indent\indent Definim următoarele funcții ce vor fi de folos:
\begin{enumerate}[(a)]
\item \textbf{Funcția lui Heaviside}: $ h : \RR \to \RR $:
  \[
    h(t) = \begin{cases}
      0, & t < 0 \\
      1, & t \geq 0
    \end{cases};
  \]
\item \textbf{Funcția sinus atenuat}: $ S_a : \RR \to \RR $:
  \[
    S_a(t) = \begin{cases}
      \frac{\sin t}{t}, & t \neq 0 \\
      1, & t = 0
    \end{cases};
  \]
\item \textbf{Funcția semn}: $ \mathrm{sgn} : \RR \to \RR $:
  \[
    \mathrm{sgn} t = \begin{cases}
      -1, & t < 0 \\
      0, & t = 0 \\
      1, & t > 0
    \end{cases}.
  \]
\end{enumerate}

Cu acestea, transformatele Fourier ale funcțiilor des întîlnite sînt redate în
tabelul din figura \ref{fig:tabel-f}.

\begin{figure}[!htbp]
  \centering
  \begin{tabular}{| c | c |}
    $ f(t) $ & $ \widehat{f}(\omega) $ \\
    \hline \hline
    $ h(t + t_0) - h(t - t_0), t_0 > 0 $ & $ 2t_0 S_a(\omega t_0) $ \\ & \\
    $ \frac{\omega_0}{\pi} S_a(\omega_0 t), \omega_0 > 0 $ & $ h(\omega + \omega_0) - h(\omega - \omega_0) $ \\ & \\
    $ h(t) - h(t_0), t_0 > 0 $ & $ t_0 S_a(\omega t_0) + \frac{1}{2i} \omega t_0^2 S_a^2 \Big( \frac{\omega t_0}{2} \Big) $\\ & \\
    $ e^{-\omega_0 t}, \omega_0 > 0 $ & $ \frac{2 \omega_0}{\omega_0^2 + \omega^2} $ \\ & \\
    $ -h(t + t_0) + 2h(t) - h(t - t_0), t_0 > 0 $ & $ \frac{1}{i} \omega t_0^2 S_a^2\Big( \frac{\omega t_0}{2} \Big) $ \\ & \\
    $ e^{-\omega_0 t} h(t), \omega_0 > 0 $ & $ \frac{\omega_0}{\omega_0^2 + \omega^2} + \frac{1}{i} \frac{\omega}{\omega_0^2 + \omega^2} $ \\ & \\
    $ e^{-\omega_0 |t|} \mathrm{sgn} t, \omega_0 > 0 $ & $ \frac{2}{i} \frac{\omega}{\omega_0^2 + \omega^2} $ \\ & \\
    $ \Big( 1 - \frac{|t|}{t_0} (h(t + t_0) - h(t - t_0)\Big), t_0 > 0 $ & $ t_0 S_a^2 \Big( \frac{\omega t_0}{2} \Big) $ \\ & \\
    $ e^{-\omega_0^2 t^2}, \omega_0 > 0 $ & $ \frac{\sqrt{\pi}{\omega_0}} e^{\frac{-\omega^2}{4 \omega_0^2}} $
  \end{tabular}
  \caption{\emph{Transformatele Fourier ale unor funcții uzuale}}
  \label{fig:tabel-f}
\end{figure}

%%%%%%%%%%%%%%%%%%%%%%%%%%%%%%%%%%%%%%%%%%%%%%%%%%%%%%%%%%%%%%%%%%%%%%

\section*{Distribuții}

\indent\indent Similar cu funcțiile de distribuție din studiul probabilităților,
introducem cîteva dintre cele mai importante funcții de distribuție, pe care le vom
trata apoi cu ajutorul transformatelor Fourier.

\begin{definition}\label{def:f-test}
  Se numește \emph{funcție test} o funcție $ \varphi : \RR \to \RR $, care este indefinit
  derivabilă și nulă în afara unui interval compact.
\end{definition}

Fie $ \kal{D} $ mulțimea tuturor funcțiilor test.

De exemplu:
\[
  \varphi_\varepsilon (x) =
  \begin{cases}
    \exp \Big( - \frac{\varepsilon^2}{\varepsilon^2 - x^2} \Big), & |x| < \varepsilon \\
    0, & |x| \geq \varepsilon
  \end{cases},
\]
care se poate defini pentru orice $ \varepsilon > 0 $ este o funcție test.

În pregătirea introducerii noțiunii de \emph{distribuție}, mai avem nevoie și de:
\begin{definition}\label{def:f-test-conv}
  Spunem că șirul de funcții test $ (\varphi_n)_n $ \emph{converge către 0} dacă acest
  șir se anulează în afara unui compact (același pentru toate) și dacă el converge
  uniform către 0 împreună cu șirul derivatelor de orice ordin.
\end{definition}

Acum, definiția principală:

\begin{definition}\label{def:distributie}
  Se numește \emph{distribuție} o aplicație liniară $ f : \kal{D} \to \RR $ cu proprietatea
  că dacă un șir de funcții test $ (\varphi_n)_n $ converge către zero, atunci
  $ \ds\lim_{n \to \infty} f(\varphi_n) = 0 $.

  Mulțimea distribuțiilor definite pe $ \kal{D} $ se va nota cu $ \kal{D}' $.
\end{definition}

\begin{remark}\label{rk:distr-sp-vec}
  Un spațiu de distribuții cu valori reale devine spațiu vectorial real, cu operațiile:
  \begin{itemize}
  \item \emph{adunarea distribuțiilor}: $ (f + g)(\varphi) = f(\varphi) + g(\varphi) $,
    pentru orice $ f, g \in \kal{D}', \varphi \in \kal{D} $;
  \item \emph{înmulțirea distribuțiilor cu scalari}: Dată distribuția $ f $ și scalarul
    $ \alpha \in \RR $, se definește distribuția $ \alpha f $ prin relația:
    \[
      (\alpha f)(\varphi) = f(\alpha \varphi) = \alpha f(\varphi), %
      \quad \forall \varphi \in \kal{D}.
    \]
  \end{itemize}
\end{remark}

Principalele distribuții pe care le vom folosi sînt \emph{distribuția Dirac} și
\emph{distribuția Heaviside}, definite mai jos.

\begin{definition}\label{def:distr-dirac}
  \emph{Distribuția Dirac} în punctul $ a $ este funcția
  \[
    \delta_a : \kal{D} \to \RR, \quad \delta_a(\varphi) = \varphi(a).
  \]

  Cazul particular $ a = 0 $ este des întîlnit, motiv pentru care vom scrie $ \delta $ în
  loc de $ \delta_0 $.
\end{definition}

\begin{definition}\label{def:distr-heaviside}
  \emph{Distribuția Heaviside} se definește prin:
  \[
    H : \kal{D} \to \RR, \quad H(\varphi) = \int_0^\infty \varphi(x) dx.
  \]
\end{definition}

O altă distribuție care se va mai întîlni este:

\begin{definition}\label{def:distr-vp}
  Distribuția $ f = VP \frac{1}{t} $ se numește \emph{valoarea principală} a lui $ \frac{1}{t} $,
  definită prin:
  \[
    f(\varphi) = v.p. \int_{-\infty}^\infty \frac{1}{t} \varphi(t) dt = %
    \lim_{\varepsilon \searrow 0} \Big( \int_{-\infty}^{-\varepsilon} \frac{\varphi(t)}{t} dt + %
    \int_\varepsilon^\infty \frac{\varphi(t)}{t} dt \Big), \quad \forall \varphi \in \kal{D}.
  \]
\end{definition}

O legătură importantă între distribuția lui Dirac și cea a lui Heaviside este următoarea:

\begin{proposition}\label{pr:dirac-h}
  Derivata distribuției lui Heaviside este distribuția lui Dirac.
\end{proposition}

\begin{proof}
  Trebuie demonstrat că $ H'(\varphi) = \delta(\varphi) $. Conform proprietății de derivare
  a distribuțiilor, $ f^{(n)}(\varphi) = (-1)^n f(\varphi^{(n)}) $, aceasta este echivalent
  cu a demonstra că $ -H(\varphi') = \delta(\varphi) $. Calculăm succesiv:
  \begin{align*}
    -H(\varphi') &= -\int_{-\infty}^\infty \varphi'(x) dx \\
                 &= - \int_0^\infty \varphi'(x) dx \\
                 &= -\varphi(x) \Big|_0^\infty \\
                 &= \varphi(0) = \delta(\varphi).
  \end{align*}
\end{proof}

Oricărei funcții \emph{local integrabile}\footnote{mărginită și integrabilă pe orice compact}
i se poate asocia o distribuție:

\begin{definition}\label{def:functie-distr}
  Fie $ u : \RR \to \RR $ o funcție local integrabilă. Pentru orice funcție test
  $ \varphi \in \kal{D} $, definim:
  \[
    \underline{u}(\varphi) = \int_{-\infty}^\infty u(x) \varphi(x) dx.
  \]
\end{definition}

Se arată simplu că $ \underline{u}(\varphi) $ este o distribuție și se numește
\emph{distribuție regulată} sau \emph{de tip funcție}.

%%%%%%%%%%%%%%%%%%%%%%%%%%%%%%%%%%%%%%%%%%%%%%%%%%%%%%%%%%%%%%%%%%%%%%

\section*{Transformarea Fourier a distribuțiilor}

\begin{definition}\label{def:f-rapid-c}
  Se numește \emph{funcție rapid crescătoare} o funcție $ f : \RR \to \CC $ de clasă
  $ \kal{C}^\infty $ pentru care toate produsele $ x^i f^{(k)} (x) $ de puteri naturale
  ale lui $ x $ și derivate ale lui $ f $ sînt funcții mărginite.

  Fie $ \kal{S} $ mulțimea acestor funcții.
\end{definition}

\begin{definition}\label{def:s-sir-c}
  Un șir $ (\psi_n)_n $ de funcții din $ \kal{S} $ se numește \emph{convergent către 0},
  notat $ \psi_n \to 0 $ în $ \kal{S} $ dacă pentru orice întregi $ j, k \geq 0 $,
  șirul de funcții $ ( x^j \psi_n^{(k)} ) $ converge uniform către 0 pe $ \RR $.
\end{definition}

Distribuțiile cărora le vom aplica transformata Fourier sînt de un tip special:

\begin{definition}\label{def:distr-temp}
  Se numește \emph{distribuție temperată} orice aplicație $ \CC $-liniară
  $ f : \kal{S} \to \CC $ astfel încît dacă $ \psi_n \to 0 $ în $ \kal{S} $,
  are loc $ \ds\lim_{n \to \infty} f(\psi_n) = 0 $.

  Se notează cu $ \kal{S}' $ spațiul tuturor distribuțiilor temperate.
\end{definition}

În fine, transformata Fourier a unei distribuții temperate este definită mai jos:

\begin{definition}\label{def:transf-f-distr}
  Fie $ f \in \kal{S}' $. Se numește \emph{transformata Fourier a distribuției $ f $}
  distribuția $ \kal{F}_f : \kal{S} \to \CC $, definită prin:
  \[
    \kal{F}_f(\psi) = f(\kal{F}(\psi)), \quad \forall \psi \in \kal{S}.
  \]
\end{definition}

De exemplu, să determinăm transformata Fourier a distribuției Dirac $ \delta $.

Calculăm succesiv:
\begin{align*}
  \kal{F}_\delta(\psi) &= \delta(\kal{F}(\psi)) \\
                       &= \kal{F}(\omega)\mid_{\omega = 0} \\
                       &= \int_{-\infty}^\infty \psi(t) e^{-i\omega t} dt \mid_{\omega = 0} \\
                       &= \int_{-\infty}^\infty \psi(t) dt \\
                       &= \underline{1}(\psi).
\end{align*}

Cum relația are loc pentru orice $ \psi \in \kal{S} $, rezultă că
$ \kal{F}_\delta = \underline{1} $.

\vspace{1cm}

Transformatele Fourier ale unor distribuții des întîlnite sînt redate în tabelul
din figura \ref{fig:transf-f-distr}.

\begin{figure}[!htbp]
  \centering
  \begin{tabular}{|c|c|}
    $ f(t) $ & $ \kal{F}_f $ \\
    \hline\hline
    $ e^{i \omega_0 t}, \omega_0 > 0 $ & $ 2 \pi \delta(\omega - \omega_0) $ \\ & \\
    $ 1 $ & $ 2 \pi \delta(\omega) $ \\ & \\
    $ \cos \omega_0 t, \omega_0 > 0 $ & $ \pi[\delta(\omega - \omega_0) + \delta(\omega + \omega_0)] $ \\ & \\
    $ \sin \omega_0 t, \omega_0 > 0 $ & $ i\pi[\delta(\omega + \omega_0) - (\omega - \omega_0)] $ \\ & \\
    $ \mathrm{sgn} t$ & $ \frac{2}{i \omega} $ \\ & \\
    $ h(t) $ & $ \pi \delta(\omega) + \frac{1}{i \omega} $ \\ & \\
    $ e^{-i \omega_0 t} h(t), \omega_0 > 0 $ & $ \pi \delta(\omega - \omega_0) + \frac{1}{i(\omega - \omega_0)} $ \\ & \\
    $ h(t) \cos \omega_0 t, \omega_0 > 0 $ & $ \frac{\pi}{2} [ \delta(\omega - \omega_0) + \delta(\omega + \omega_0)] + \frac{i \omega}{\omega_0^2 - \omega^2} $ \\ & \\
    $ h(t) \sin \omega_0 t, \omega_0 > 0 $ & $ \frac{\pi}{2i} [ \delta(\omega - \omega_0) + \delta(\omega + \omega_0)] + \frac{\omega_0}{\omega_0^2 - \omega^2} $ \\ & \\
    $ t^n $ & $ 2 \pi i^n \frac{d^n \delta(\omega)}{d \omega^n} $ \\
    \hline
  \end{tabular}
  \caption{\emph{Transformatele Fourier ale unor distribuții uzuale}}
  \label{fig:transf-f-distr}
\end{figure}

%%%%%%%%%%%%%%%%%%%%%%%%%%%%%%%%%%%%%%%%%%%%%%%%%%%%%%%%%%%%%%%%%%%%%%

\section*{Exerciții}

\indent\indent 1. Găsiți transformata Fourier a funcției:
\[
  f(x) =
  \begin{cases}
    1, & |x| < 1 \\
    0, & \text{în rest}
  \end{cases}.
\]

\emph{Soluție:} Folosim definiția


  
\end{document}